% PAGES: 246-246

respectively, where $N(z)$ is the cumulative distribution of the standard normal
variable, i.e.,
\[
N(z) \;=\; \frac{1}{\sqrt{2\pi}} \int_{-\infty}^{z} e^{-\frac{x^2}{2}} \, dx,
\]
and
\begin{equation}
d_1 \;=\; \frac{\ln\left(\frac{S}{K}\right) + \left(r-q+\frac{\sigma^2}{2}\right)T}
{\sigma\sqrt{T}}; \quad d_2 \;=\; d_1 - \sigma\sqrt{T}.
\end{equation}

For any plain vanilla European option, the option price $C_m$ or $P_m$, the maturity
$T$, the strike $K$, and the spot price $S$ of the underlying asset are known. However,
a continuous dividend yield $q$ for the underlying asset is very rarely quoted in the
markets and the interest rate $r$ could be chosen from several different discount
curves. As seen below,\footnote{This method was implemented in 2010 in Bloomberg
terminals, providing a tenfold improvement in the accuracy of implied volatility
calculations over the prior method.} the least squares method and Put--Call parity can
be used to overcome these issues and find the implied volatility using Newton's method.

As an example of how implied volatilities are computed in practice, consider the
snapshot from Table 8.2 of the mid prices\footnote{The mid price of an option is the
average of the bid price and ask price of the option.} on March 9, 2012, of the S\&P
500 options (ticker symbol SPX)\footnote{More information on SPX options can be found
on the Chicago Board Options Exchange (CBOE); see
http://www.cboe.com/products/indexopts/spx\_spec.aspx} maturing on December 22, 2012.
These options are European options and therefore we can use the Black--Scholes
framework. Although not needed for this method, the corresponding spot price of the
index was $1,370$.

\begin{table}[H]
\centering
\caption{Dec 2012 SPX option prices on 3/9/2012}
\begin{tabular}{|c|c||c|c|}
\hline
Call Strike & Price & Put Strike & Price \\
\hline
C1175 & 225.40 & P1175 & 46.60 \\
\hline
C1200 & 205.55 & P1200 & 51.55 \\
\hline
C1225 & 186.20 & P1225 & 57.15 \\
\hline
C1250 & 167.50 & P1250 & 63.30 \\
\hline
C1275 & 149.15 & P1275 & 70.15 \\
\hline
C1300 & 131.70 & P1300 & 77.70 \\
\hline
C1325 & 115.25 & P1325 & 86.20 \\
\hline
C1350 & 99.55 & P1350 & 95.30 \\
\hline
C1375 & 84.90 & P1375 & 105.30 \\
\hline
C1400 & 71.10 & P1400 & 116.55 \\
\hline
C1425 & 58.70 & P1425 & 129.00 \\
\hline
C1450 & 47.25 & P1450 & 143.20 \\
\hline
C1500 & 29.25 & P1500 & 173.95 \\
\hline
C1550 & 15.80 & P1550 & 210.80 \\
\hline
C1575 & 11.10 & P1575 & 230.90 \\
\hline
C1600 & 7.90 & P1600 & 252.40 \\
\hline
\end{tabular}
\end{table}

Recall from Section 10.3 that the Put--Call parity states that taking a long position
in a European call option and a short position in a European put option with the same
strike $K$ and maturity $T$ is equivalent to taking a long position in a forward
contract with delivery price $K$ and maturity $T$, and therefore the following
